\section*{अध्याय \ref{ch:g8:the-oceans} --- महासागर}

\begin{solution}{exo:g8:the-oceans:1}
\omterm{def:g8:the-oceans:seafloor}{महाद्वीपीय मग्नतट}, \omterm{def:g8:the-oceans:seafloor}{महाद्वीपीय ढाल}, \omterm{def:g8:the-oceans:seafloor}{वितलीय मैदान}, \omterm{def:g8:the-oceans:seafloor}{महासागरीय कटक}।
\end{solution}

\begin{solution}{exo:g8:the-oceans:2}
(क) ग़लत: लगभग \qty{3700}{m}, यानी \qty{3.7}{km}। (ख) सही। (ग) ग़लत: वाष्पीकरण
पानी हटाता है और लवण छोड़ देता है, इसलिए पानी ज़्यादा खारा हो जाता है।
(घ) सही।
\end{solution}

\begin{solution}{exo:g8:the-oceans:3}
(क) \omterm{def:g8:the-oceans:seafloor}{महासागरीय कटक}। (ख) \omterm{def:g8:the-oceans:salinity}{लवणता}। (ग) \omterm{def:g8:the-oceans:current}{वलय}। (घ) \omterm{def:g8:the-oceans:seafloor}{महासागरीय खाई}।
\end{solution}

\begin{solution}{exo:g8:the-oceans:4}
(क) वामावर्त। (ख) गरम: गल्फ़ स्ट्रीम; ठंडी: कैनरी धारा। (ग) \omterm{def:g8:winds-global-circulation:global-winds}{व्यापारिक
पवनें}।
\end{solution}

\begin{solution}{exo:g8:the-oceans:5}
(क) $35 \div 1000 \times 100 = 3.5~\%$। (ख) $20 \times 35 = \qty{700}{g}$।
(ग) $1000 \div 35 \approx 28.6$: लगभग \qty{29}{kg}।
\end{solution}

\begin{solution}{exo:g8:the-oceans:6}
(क) $1500 \times 6 = \qty{9000}{m}$। (ख) ध्वनि तली तक नीचे जाती है और वापस ऊपर आती
है: वह गहराई को दो बार तय करती है। (ग) $9000 \div 2 =
\qty{4500}{m}$। (घ) एक \omterm{def:g8:the-oceans:seafloor}{वितलीय
मैदान} (खाई ज़्यादा गहरी होती, कटक या मग्नतट कम
गहरा)।
\end{solution}

\begin{solution}{exo:g8:the-oceans:7}
(क) लगभग \qty{3500}{m} और \qty{5000}{m} के बीच। (ख) प्रोफ़ाइल के साथ लगभग
\qty{4400}{km} पर, लगभग \qty{2300}{m} गहरा। (ग) लगभग
$4500 - 2300 = \qty{2200}{m}$ (1200 और \qty{2700}{m} के बीच, इस पर निर्भर कि कौन-सा
\omterm{def:g1:landscapes-around-us:landform}{मैदान} चुना गया)। (घ) \omterm{def:g2:water-on-earth:ocean}{महासागर} कुछ किलोमीटर गहरा पर हज़ारों किलोमीटर चौड़ा है: एक ही
पैमाने पर बनाने से पूरा प्रोफ़ाइल एक रेखा से भी पतला होता।
\end{solution}

\begin{solution}{exo:g8:the-oceans:8}
\omterm{def:g2:water-on-earth:ocean}{समुद्र} A ज़्यादा खारा है: सूर्य और सूखी हवा तेज़ \emph{वाष्पीकरण} करते हैं, पानी चला
जाता है, \emph{लवण पीछे रह जाते हैं}, और लगभग कोई वर्षा या \omterm{def:g1:landscapes-around-us:water}{नदी} \emph{\omterm{def:g2:water-on-earth:salt-fresh}{मीठा पानी}} वापस
नहीं लाती। \omterm{def:g2:water-on-earth:ocean}{समुद्र} B कम खारा है: उसकी \omterm{def:g1:landscapes-around-us:water}{नदियाँ} और उसकी वर्षा बहुत \emph{\omterm{def:g2:water-on-earth:salt-fresh}{मीठा पानी}} लाती
हैं, जो लवणों को पतला कर देता
है।
\end{solution}

\begin{solution}{exo:g8:the-oceans:9}
(क) गल्फ़ स्ट्रीम, \omterm{def:g1:landscapes-around-us:water}{तट} के साथ-साथ उत्तर की ओर, फिर उत्तर-पूर्व की ओर। (ख) पूर्व की ओर
(उत्तर-पूर्व की ओर), उत्तर अटलांटिक धारा के साथ। (ग) उत्तर-पश्चिमी अफ़्रीका (और स्पेन
और पुर्तगाल) के \omterm{def:g1:landscapes-around-us:water}{तट} के साथ-साथ, कैनरी धारा के साथ। (घ)
ठंडी।
\end{solution}

\begin{solution}{exo:g8:the-oceans:10}
(क) $6.4 \times 24 = 153.6$: लगभग \qty{154}{km}। (ख) $153.6 \times 7
\approx \qty{1075}{km}$।
(ग) गल्फ़ स्ट्रीम पैदल चलने वाले से थोड़ी तेज़ है, पर वह दिन-रात कभी नहीं
रुकती।
\end{solution}

\begin{solution}{exo:g8:the-oceans:11}
(क) \qty{3.7}{km} और \qty{11}{km}। (ख) $11000 \div 3700 \approx 3.0$
गुना। (ग) ध्वनि $2 \times 11000 = \qty{22000}{m}$ चलती है;
$22000
\div 1500 \approx 14.7$: लगभग \qty{15}{s}। (घ) $1000 = 10^3$ वर्ष।
\end{solution}

\begin{solution}{exo:g8:the-oceans:12}
(क) \omterm{def:g3:continents-and-oceans:equator}{भूमध्य रेखा} के पास \omterm{def:g8:winds-global-circulation:global-winds}{व्यापारिक पवनें} गरम उष्णकटिबंधीय पानी को पश्चिम की ओर धकेलती हैं।
अटलांटिक के पश्चिमी किनारे पर वह गल्फ़ स्ट्रीम के साथ, जो एक गरम धारा है, फ़्लोरिडा और
उत्तर अमेरिका के साथ-साथ उत्तर की ओर मुड़ता है। फिर \omterm{def:g8:winds-global-circulation:global-winds}{पछुआ पवनें} उसे उत्तर-पूर्व की ओर
अटलांटिक के पार, नॉर्वे के सामने के \omterm{def:g2:water-on-earth:ocean}{समुद्रों} तक, धकेलती हैं। (ख) वहाँ वह अपनी ऊष्मा
ठंडी हवा को देता है: वह ठंडा होता है, और समुद्री बर्फ़ बनने पर ज़्यादा खारा हो जाता है,
क्योंकि लवण पानी में रह जाता है। ठंडा, \omterm{def:g2:water-on-earth:salt-fresh}{खारा पानी} घना होता है, इसलिए वह तली तक डूब जाता
है। (ग) वह \omterm{def:g2:water-on-earth:ocean}{महासागर} की तली के साथ-साथ दक्षिण की ओर, अंटार्कटिका की ओर, फिर हिंद और
प्रशांत \omterm{def:g2:water-on-earth:ocean}{महासागरों} में रेंगता है, जहाँ वह धीरे-धीरे फिर ऊपर उठता है।
\end{solution}

\begin{solution}{pb:g8:the-oceans:1}
\textbf{1.} अटलांटिक \omterm{def:g2:water-on-earth:ocean}{महासागर} में, \omterm{def:g3:continents-and-oceans:equator}{उत्तरी गोलार्ध} में।

\textbf{2.} दक्षिणावर्त; शुरू में प्लव उत्तर की ओर ले जाया जाता है।

\textbf{3.} हिस्सा A (गल्फ़ स्ट्रीम) गरम है: वह \omterm{def:g5:weather-and-climate:tropics}{उष्णकटिबंध} से आता है। हिस्सा C
(कैनरी धारा) ठंडा है: वह ऊँचे \omterm{def:g6:latitude-longitude:latitude}{अक्षांशों} से, \omterm{def:g3:continents-and-oceans:equator}{भूमध्य रेखा} की ओर,
आता है।

\textbf{4.} हिस्सा B: \omterm{def:g8:winds-global-circulation:global-winds}{पछुआ पवनें}। हिस्सा D: \omterm{def:g8:winds-global-circulation:global-winds}{व्यापारिक पवनें}।

\textbf{5.} $3800 \div 100 = 38$ दिन।

\textbf{6.} B: $2100 \div 20 = 105$ दिन। C: $3100 \div 20 = 155$ दिन।
D: $6100 \div 20 = 305$ दिन।

\textbf{7.} $3800 + 2100 + 3100 + 6100 = \qty{15100}{km}$।

\textbf{8.} गल्फ़ स्ट्रीम \omterm{def:g8:the-oceans:current}{वलय} के पश्चिमी किनारे की संकरी, तेज़ धारा है; बाकी किनारों
की धाराएँ चौड़ी और धीमी हैं।

\textbf{9.} $38 + 105 + 155 + 305 = 603$ दिन।

\textbf{10.} $603 \div 30 = 20.1$: लगभग \textbf{20 महीने}।

\textbf{11.} हिस्से A, B और C के बाद $38 + 105 + 155 = 298$ दिन बीत चुके होते हैं;
365 दिन हिस्से D (उत्तर विषुवतीय धारा) के दौरान पूरे होते हैं, उसमें
$365 -
298 = 67$ दिन बाद।

\textbf{12.} $603 - 365 = 238$ दिन और, यानी $238 \div 30 \approx 8$ महीने
और।
\end{solution}
